\documentclass[10pt]{article}
\usepackage[margin=0.95in]{geometry}
\usepackage{amsmath,amssymb,amsthm,microtype}
\usepackage[colorlinks=true,urlcolor=blue!50!black]{hyperref}
\usepackage{xcolor}
\newtheorem{theorem}{Theorem}
\newtheorem{lemma}{Lemma}
\newcommand{\floor}[1]{\left\lfloor#1\right\rfloor}
\title{Excluding All Irrational Mechanical\\Collatz Itineraries in Lean}
\author{Collatz proof project}
\date{September 5, 2026}
\begin{document}
\maketitle
\begin{abstract}
We connect the repository's parity factor-complexity obstruction to its
mechanical-word definition. Every mechanical parity itinerary has at most
$L+1$ distinct length-$L$ factors, hence its natural Collatz orbit is bounded.
A return of period $q$ forces $q\alpha$ to be an integer. Consequently no
natural seed, including any shifted seed, has a mechanical itinerary of
irrational slope, for any real intercept. The proof is checked in Lean
without added mathematical axioms. It closes a formal coverage gap between
previous slope-specific results and a full irrational mechanical class
exclusion. It does not prove Collatz; no mathematical priority is claimed.
\end{abstract}

\section{Definitions and statement}
Let $T(n)=n/2$ for even $n$ and $T(n)=(3n+1)/2$ for odd $n$.
Write $x_t=T^t(N)$ and let $p_N(L)$ count all distinct factors
$(x_t\bmod2,\ldots,x_{t+L-1}\bmod2)$ for $t\ge0$.
For real $\alpha,\rho$, define
\[
w_t=\floor{(t+1)\alpha+\rho}-\floor{t\alpha+\rho}.
\]
The Lean predicate \texttt{HasItin} equates these integer letters with
the actual parity values for every natural time.
\begin{theorem}
For every irrational real $\alpha$, real $\rho$, and natural $N$, the
parity itinerary of $N$ is not $w$. The same holds for every tail $T^s(N)$.
\end{theorem}
No critical-density, continued-fraction, or positive-seed hypothesis is
needed. The complexity implication below even allows arbitrary real slopes;
the itinerary hypothesis itself enforces binary letters.

\section{The ordered-threshold argument}
Normalize a starting time $t$ to
$\beta_t=\{t\alpha+\rho\}\in[0,1)$.
Subtracting the integer part shows that the length-$L$ factor is determined
by $\floor{i\alpha+\beta_t}$ for $0\le i\le L$.
For $1\le i\le L$, put
\[
d_i(\beta)=\floor{i\alpha+\beta}-\floor{i\alpha},\qquad
R(\beta)=\sum_{i=1}^L d_i(\beta).
\]
Each $d_i$ is nondecreasing and takes values in $\{0,1\}$, so
$0\le R(\beta)\le L$. If two intercepts have the same rank, order them.
All coordinate differences then have the same sign and sum to zero;
thus every coordinate agrees. Their successive floor differences agree
as well. Choose one realizing start for each factor and assign its rank.
This is an injection into $\{0,\ldots,L\}$, proving
\[
p_N(L)\le L+1.
\]
The finite set of globally realized factors and the choice of representatives
are noncomputable definitions, not an orbit-enumeration algorithm.

\section{Boundedness and the irrational contradiction}
The previously verified factor floor states that an unbounded orbit satisfies
\[
5q+1\le p_N(3q+K)\quad\text{whenever }N+1\le2^K.
\]
Combining it with $p_N(L)\le L+1$ contradicts sufficiently large $q$.
Thus a mechanical itinerary forces boundedness. Pigeonhole gives repeated
states, and shifting to the first repeated state gives a return of some
positive period $q$.

Let $p$ be the odd-step count of that period. Determinism and additivity
imply that $k$ periods contain $kp$ odd steps. Mechanical telescoping yields
\[
kp=\floor{kq\alpha+\rho'}-\floor{\rho'},\qquad
-1<k(q\alpha-p)<1\quad(k\ge0).
\]
If $q\alpha-p$ were nonzero, the Archimedean property would violate one of
these bounds for a sufficiently large natural $k$. Therefore $q\alpha=p$,
contradicting irrationality because $q>0$.

\section{Formal declarations and verification}
All new declarations are in \texttt{Collatz/MechanicalComplexity.lean}:
\begin{itemize}
\item \texttt{mechanical\_complexity\_le}: the global factor upper bound;
\item \texttt{bounded\_of\_mechanical}: boundedness for any mechanical itinerary;
\item \texttt{mechanical\_return\_slope}: integral slope times return period;
\item \texttt{no\_irrational\_mechanical}: the full irrational-slope exclusion;
\item \texttt{no\_eventual\_irrational\_mechanical}: every tail is excluded.
\end{itemize}
The module builds under the repository's Lean toolchain. These five
statements are included in the selected axiom audit; their dependencies
are confined to the standard foundations \texttt{propext},
\texttt{Classical.choice}, and \texttt{Quot.sound}. No external Sturmian
equivalence or literature axiom is imported. This is a Lean build and
axiom inspection, not an independent kernel replay.

\section{Context and remaining problem}
Dubickas proved a stronger asymptotic complexity lower bound for escaping
Collatz trajectories in Theorem~5 of \emph{On integer sequences generated
by linear maps}, Glasgow Mathematical Journal 51 (2009), 243--252,
\href{https://doi.org/10.1017/S0017089508004655}{doi:10.1017/S0017089508004655}.
Our contribution here is the explicit formal bridge to mechanical words,
using the repository's weaker finite rational complexity floor.

This theorem subsumes the earlier irrational slope-family exclusions at
the level of conclusions, while preserving their different quantitative
proofs. It gives neither a complexity upper bound for arbitrary Collatz
orbits nor a classification of all possible counterexamples. For rational
mechanical slopes it proves boundedness but does not identify the eventual
cycle. General escape and nontrivial cycles remain unresolved.
\end{document}
